\chapter{Unsupervised Learning and Representation}\label{chap:unsupervised}
\courseweek{Week 6.} Unsupervised methods organize data without a labelled target. They are useful for exploration, anomaly investigation, compression, and finding candidate subpopulations. They do not automatically create ground truth.

\begin{learningbox}
\begin{itemize}
  \item Explain the difference between clustering and classification.
  \item Standardize features appropriately before distance-based analysis.
  \item Use a low-dimensional view as an exploratory tool rather than a final conclusion.
  \item State how a discovered cluster would be checked against domain knowledge.
\end{itemize}
\end{learningbox}

\section{Clustering as a hypothesis generator}
Methods such as $k$-means partition observations by geometric similarity. The algorithm returns groups even when no meaningful groups exist. Therefore a cluster label should be treated as a hypothesis to investigate with process knowledge, additional measurements, and sensitivity analysis.

\section{Representation matters}
Distances depend on feature scale and representation. A temperature difference of 10 degrees and a pressure difference of 10 kPa are not inherently comparable. Document scaling, missing-value handling, and the variables intentionally excluded from the analysis.

\begin{labbox}
\textbf{Week 6:} standardize a multivariate dataset, compare two values of $k$, and inspect the stability of the result. Identify one observation that would require domain review rather than automatic labelling.
\end{labbox}
